\section{What Drives Base-Model Choice}
\label{sec:results}

\begin{table}[t]
\centering
\caption{Shapley decomposition of the log-likelihood gain over uniform attachment (main specification: $(d{+}1)$ PPS candidate sets, first adoptions). In-sample shares with 95\% multiplier-bootstrap intervals (50 replicates, authors resampled); holdout = mean per-choice log-likelihood on the last 10\% of choices in time.}
\label{tab:shapley}
\small
\begin{tabular}{lrcr}
\toprule
Group & Share (\%) & 95\% CI & Holdout (\%) \\
\midrule
Popularity & 91.6 & [91.1, 92.0] & 91.8 \\
General quality & 0.006 & [0.003, 0.010] & 0.001 \\
Demand match & 5.2 & [5.0, 5.4] & 4.6 \\
Inertia & 3.2 & [2.7, 3.6] & 3.6 \\
\midrule
\multicolumn{4}{l}{Holdout gain over uniform: 6.62 nats/choice; top-1 / top-10 accuracy 14.2\% / 39.8\%} \\
\bottomrule
\end{tabular}
\end{table}


\subsection{Popularity explains almost all of what can be explained}
\Cref{tab:shapley} reports the decomposition for the main specification. Popularity accounts for
\shPopularity\% of the log-likelihood gain over uniform attachment (95\% CI \shPopularityLo--\shPopularityHi),
demand match for \shMatch\% (\shMatchLo--\shMatchHi), inertia for \shInertia\% (\shInertiaLo--\shInertiaHi) and
general quality for less than 0.01\%. The holdout shares are nearly identical (\shHoldPopularity\%, \shHoldMatch\%,
\shHoldInertia\%), so the ranking is not an artifact of in-sample fit. Every group is jointly significant (likelihood
ratio for dropping popularity \lrPop, demand match \lrMatch, inertia \lrInertia, quality \lrQual{} on 1 d.f.); the
question is not whether developers respond to fit but how much of their choices fit accounts for. Out of sample the
model gains \holdGain{} nats per choice over uniform attachment and ranks the chosen parent first among the
$\approx$300 sampled candidates in \holdTopOne\% of holdout choices (top-10: \holdTopTen\%).

\paragraph{Developers value fit, but fit rarely separates candidates.}
The small share of demand match is not because developers ignore it. In utility terms (\cref{tab:coef}), a
candidate that matches the task of the author's earlier parents is worth as much as a \mrsTask-fold increase in
$1+d$ (95\% CI \mrsTaskLo--\mrsTaskHi); one standard deviation of recent adoptions is worth an \mrsRec-fold increase,
a major-organization publisher \mrsOrg-fold. Per unit, fit is valued highly. It explains little of the variation in
choices because the two kinds of signal vary very differently across candidates: degree spans five orders of
magnitude within a candidate set, whereas the match signals are coarse, and many candidates fit equally well. Among
first adoptions with a demand signal, on average \tieMean{} sampled candidates tie at the highest match index, and
what separates them is popularity. Fit narrows the field; popularity picks the winner.

\paragraph{Quality.}
Model-index benchmark scores are reported for well under 1\% of candidates, and their coefficient is small and, if
anything, negative. This could mean that developers do not respond to quality, or only that quality is rarely
measured in the card. The next subsection bounds the second possibility.

\begin{table}[t]
\centering
\caption{How large could fit be? Fit = general quality + demand match. Hub FE: 500 hub-specific intercepts (top-500 by degree at the holdout cut). External quality explains 7.3\% of the FE variance among the 118 hubs covered by the Open LLM Leaderboard v2 (noise-corrected, with redistributor and size controls).}
\label{tab:bounds}
\small
\begin{tabular}{lrr}
\toprule
Specification & Fit (\%) & Popularity (\%) \\
\midrule
No FE; quality = benchmark z-score & 5.2 & 91.6 \\
Hub FE counted as quality (in-sample) & 46.0 & 51.0 \\
Hub FE counted as quality (holdout) & 36.2 & 60.3 \\
Externally explained part of FE (in-sample) & 7.0 & 89.9 \\
Externally explained part of FE (holdout) & 6.3 & 90.2 \\
\bottomrule
\end{tabular}
\end{table}


\subsection{How large could fit be?}
\paragraph{An upper bound from hub fixed effects.}
Adding fixed effects for the 500 most adopted parents, and counting them entirely as quality, is the most generous
reading of quality: any fixed attribute of a hub, measured or not, would be credited to it. Under this reading fit
would account for \feInQual\% of the in-sample gain and \feHoldQual\% of the holdout gain (\cref{tab:bounds}). Two
observations show that the bound is far from tight. First, the fixed effects add almost nothing out of sample:
the holdout gain moves from \holdGainNoFE{} to \holdGainFE{} nats per choice. In-sample, the fixed effects take
over credit from the popularity features they are collinear with (the popularity share falls from \shPopularity\%
to 51.0\%), but they predict future choices no better than the popularity signals already do. Second, external
quality explains little of them: among the \extHubs{} hubs covered by the Open LLM Leaderboard v2, the
noise-corrected share of fixed-effect variance explained by the leaderboard score (with size and re-distributor
controls) is \extQ\%. Splitting the fixed-effect share accordingly, fit accounts for \extFitIn\% in-sample and
\extFitHold\% on the holdout. Whatever makes a hub attractive beyond its observed popularity, it is mostly not what
benchmarks measure.\todo{Caveat: leaderboard coverage is 118 of 500 hubs; discuss representativeness.}

\paragraph{Is the demand-match share developer-specific demand?}
Some of what looks like match could be common trends: if everyone is fine-tuning instruction-following text models
in a given month, any developer's history will ``match'' the models that are popular that month. We replace each
developer's demand with that of a donor, another author's first adoption in the same month and modality, and
recompute the match features. On its own the donor's demand still accounts for \placeboAlone\% of the gain, so common
trends are real. But when own and donor match enter the same model, own match keeps \jointOwn\% of the gain and the
donor \jointDonor\%, and the own-match coefficients barely move (task history 1.93 vs.\ 1.90). The donor's demand
competes with popularity, not with the developer's own demand: about \jointOwn\% of the gain is fit to a
developer-specific need.

\paragraph{Popularity without a quality difference.}
The cleanest evidence that popularity signals act on their own comes from \idN{} first adoptions in which two copies
of the same weights (identical weight hashes, published by different accounts) are both in the candidate set and
one is chosen. Quality is identical by construction. The chosen copy is the one with more recent adoptions
(coefficient \idRec, SE \idRecSe), not the one with more adoptions overall (\idDeg, SE \idDegSe; the chosen copy has
the higher degree in \idHigher\% of cases): among exact substitutes, developers follow current momentum.

\begin{table*}[t]
\centering
\caption{The decomposition under three candidate-set designs. Same choices, features and estimator; only the inclusion probabilities $\pi$ differ. Coefficients with bootstrap SEs. Under uniform sampling almost every sampled competitor is a never-used model, so top-1 accuracy is high but the data contain few comparisons among popular models.}
\label{tab:q}
\footnotesize
\setlength{\tabcolsep}{4pt}
\begin{tabular}{lccc}
\toprule
 & $(d{+}1)$ PPS & Stratified & Uniform \\
\midrule
Popularity share (\%) & 91.6 [91.1, 92.0] & 92.2 [92.0, 92.5] & 88.3 [87.8, 88.6] \\
General quality share (\%) & 0.006 [0.003, 0.010] & 0.006 [0.003, 0.010] & 0.009 [0.005, 0.015] \\
Demand match share (\%) & 5.2 [5.0, 5.4] & 5.6 [5.4, 5.9] & 8.4 [8.1, 8.7] \\
Inertia share (\%) & 3.2 [2.7, 3.6] & 2.1 [1.9, 2.4] & 3.4 [3.1, 3.8] \\
\midrule
$\log d$ & 0.34 (0.01) & 0.34 (0.01) & 0.44 (0.03) \\
$\mathbf{1}[d>0]$ & 4.64 (0.09) & 4.71 (0.09) & 3.25 (0.10) \\
major org. & 0.65 (0.03) & 0.68 (0.03) & 1.48 (0.05) \\
language history & 0.25 (0.03) & 0.24 (0.03) & 0.95 (0.06) \\
task history & 1.90 (0.04) & 1.88 (0.04) & 2.38 (0.08) \\
card-text sim. & 1.69 (1.64) & 4.36 (0.90) & 0.74 (2.66) \\
\midrule
Holdout top-1 acc. (\%) & 14.2 & 13.5 & 93.1 \\
Bootstrap replicates (shares / coef.) & 50 / 200 & 20 / 80 & 20 / 80 \\
\bottomrule
\end{tabular}
\end{table*}


\subsection{Robustness to the candidate-set design}
\label{sec:robustq}
\Cref{tab:q} re-estimates the model with the same choices and features under three inclusion schemes: the main
$(d+1)$ PPS design, a stratified design (all top-500 hubs plus random positive-degree and zero-degree models) and
uniform sampling. PPS and stratified sets give the same answer: shares within one percentage point and coefficients
within one standard error. Both contain the popular competitors, and the $-\log\pi$ correction handles the rest.

Uniform sets give different weights. Almost every sampled competitor is a never-used model, so the chosen parent is
easy to find (top-1 holdout accuracy \topOneUnif\%), but the choice is then mostly a comparison between ``known'' and
``unknown'' models. On this margin the weights on organization brand, language history and the library ecosystem
roughly double, $\1[d>0]$ falls from 4.64 to 3.25, and the match share rises to 8.4\%. Since all three designs
recover the true parameters on synthetic data (\cref{sec:validation}), the difference is not an estimation
artifact: a single logit cannot describe both the known-versus-unknown margin and the choice among popular models
with the same weights. We take the PPS design as the main specification because it represents the comparisons
developers actually make, among the incumbents of their modality.

\paragraph{Other checks.}
Demand match can only register where a demand signal exists; \sigShare\% of first adoptions have one (author
history, dataset profile or card text). On these choices demand match accounts for \sigMatch\% of the gain
(\sigMatchRw\% after reweighting them to the composition of all choices by author experience and organization), and
popularity for \sigPop\%. Full-data estimates agree with an earlier analysis of a 10\% sample: \valIn{} of \valN{}
coefficients lie in the 10\% confidence interval, and the exceptions are features redefined since
(\cref{app:validation}).
